\documentclass[../DoAn.tex]{subfiles}
\begin{document}

\section{Mở đầu chương}

Chương này trình bày cơ sở lý thuyết phục vụ việc xây dựng trợ lý học vụ
\textit{Retrieval-Augmented Generation} (RAG) trong đồ án. Khác với một chatbot
chỉ sinh văn bản từ mô hình ngôn ngữ, hệ thống RAG phải kết hợp nhiều lớp kiến
thức: xử lý tài liệu, chia đoạn và metadata, biểu diễn ngữ nghĩa, tìm kiếm từ
khóa, hợp nhất kết quả truy hồi, xếp hạng lại, hiểu truy vấn hội thoại, sinh câu
trả lời có căn cứ và đánh giá chất lượng.

Nội dung của chương được tổ chức theo các nền tảng đang được sử dụng trong hệ
thống \texttt{RAG\_v2}: đặc thù dữ liệu học vụ; sự khác biệt giữa tìm kiếm truyền
thống và RAG; mô hình ngôn ngữ lớn; embedding và vector database; BM25 và hybrid
retrieval; metadata-aware retrieval; query routing, reflection và decomposition;
reranking; chunking tài liệu; Agentic RAG; web fallback; truy vết và đánh giá.
Các phần triển khai cụ thể của từng module sẽ được trình bày ở các chương sau.

\section{Bối cảnh bài toán hỏi đáp học vụ}

\subsection{Đặc thù nguồn tri thức học vụ}

Nguồn tri thức học vụ có tính dị thể cao. Một câu hỏi của sinh viên có thể liên
quan đến quy định ổn định trong nhiều năm, chương trình đào tạo của một ngành,
kế hoạch theo học kỳ hoặc thủ tục hỗ trợ sinh viên. Vì vậy, dữ liệu không nên
được xem là một tập văn bản đồng nhất. Trong hệ thống hiện tại, các nguồn tri
thức chính được tổ chức theo bốn miền:

\begin{table}[h]
\centering
\begin{tabular}{|l|p{4.2cm}|p{6.3cm}|}
\hline
\textbf{Collection} & \textbf{Nội dung chính} & \textbf{Đặc điểm truy hồi} \\
\hline
\texttt{ctdt} & Chương trình đào tạo, học phần, tín chỉ, học phần tiên quyết &
Phụ thuộc mạnh vào ngành, mã ngành và cấu trúc bảng học phần. \\
\hline
\texttt{quydinh} & Quy chế, học bổng, tốt nghiệp, ngoại ngữ, phạm vi áp dụng &
Cần chú ý khóa áp dụng, ngành áp dụng, hiệu lực văn bản và tham chiếu điều khoản. \\
\hline
\texttt{kehoach} & Kế hoạch, lịch, thông báo, thời hạn đăng ký &
Có tính thời điểm; ngày đăng bài và yêu cầu ``mới nhất'' ảnh hưởng đến kết quả. \\
\hline
\texttt{stsv} & Sổ tay, biểu mẫu, thủ tục và dịch vụ hỗ trợ sinh viên &
Nội dung thường mang tính hướng dẫn, cấu trúc khác văn bản quy định. \\
\hline
\end{tabular}
\caption{Các miền tri thức chính của trợ lý học vụ}
\end{table}

Sự phân miền dữ liệu giúp giảm nhiễu khi tìm kiếm. Chẳng hạn, câu hỏi về học phần
tiên quyết thường cần ưu tiên \texttt{ctdt}, trong khi câu hỏi về lịch đăng ký
học tập gần nhất cần ưu tiên \texttt{kehoach}. Nếu chỉ gom tất cả tài liệu vào
một kho chung, các đoạn văn có từ khóa giống nhau nhưng khác mục đích có thể cạnh
tranh sai trong danh sách kết quả.

\subsection{Đặc điểm truy vấn của sinh viên}

Truy vấn học vụ thường có các đặc điểm sau:

\begin{itemize}
    \item người dùng hỏi bằng tiếng Việt tự nhiên, không nhất thiết dùng đúng từ
    khóa xuất hiện trong tài liệu;
    \item nhiều truy vấn chứa thực thể có cấu trúc như mã ngành, mã học phần,
    khóa sinh viên, học kỳ hoặc năm học;
    \item câu hỏi nối tiếp có thể dùng đại từ hoặc cụm tham chiếu như ``môn đó'',
    ``ngành của em'', ``so sánh với chương trình kia'';
    \item một truy vấn có thể cần nhiều nguồn, ví dụ vừa hỏi chương trình học vừa
    hỏi điều kiện áp dụng của quy định;
    \item câu hỏi về kế hoạch hoặc thời hạn có thể yêu cầu dữ liệu mới hơn dữ liệu
    đã lập chỉ mục.
\end{itemize}

Các đặc điểm này tạo ra yêu cầu vừa truy hồi theo ngữ nghĩa, vừa giữ được các từ
khóa chính xác, đồng thời khai thác metadata và lịch sử hội thoại một cách có
kiểm soát.

\subsection{Yêu cầu đối với một hệ thống trả lời có căn cứ}

Một trợ lý học vụ cần ưu tiên tính kiểm chứng thay vì chỉ sinh câu trả lời trôi
chảy. Với miền tri thức này, một hệ thống hỏi đáp cần:

\begin{itemize}
    \item truy hồi các đoạn tài liệu liên quan trước khi sinh câu trả lời;
    \item giữ thông tin nguồn, collection và metadata của đoạn truy hồi;
    \item phân biệt văn bản hiện hành với văn bản đã bị thay thế khi có lineage;
    \item xử lý dẫn chiếu trong văn bản quy định như điều và khoản;
    \item tránh đưa thông tin cá nhân hoặc phạm vi học kỳ cũ vào truy vấn chung
    nếu câu hỏi hiện tại không yêu cầu;
    \item có quy trình đánh giá riêng cho factual retrieval hiện hành và cho chất
    lượng hội thoại lịch sử.
\end{itemize}

\section{Từ tìm kiếm truyền thống đến Retrieval-Augmented Generation}

\subsection{Hỏi đáp dựa trên từ khóa}

Các hệ thống tra cứu truyền thống thường bắt đầu từ keyword matching, TF--IDF
hoặc BM25. Nhóm phương pháp này phù hợp với truy vấn chứa cụm từ chính xác, mã
học phần, tên biểu mẫu hoặc số điều khoản. Tuy nhiên, nếu người dùng diễn đạt
khác tài liệu, tìm kiếm từ khóa có thể bỏ sót đoạn văn phù hợp.

Ví dụ, câu hỏi ``học phần nào phải học trước môn này'' có thể không chứa cụm từ
``học phần tiên quyết'' xuất hiện trong tài liệu. Trong trường hợp đó, tìm kiếm
ngữ nghĩa giúp thu hẹp khoảng cách diễn đạt giữa người dùng và văn bản.

\subsection{Mô hình ngôn ngữ lớn và giới hạn của sinh văn bản thuần túy}

Mô hình ngôn ngữ lớn (Large Language Model, LLM) có khả năng tổng hợp ngữ cảnh,
diễn đạt câu trả lời tự nhiên và xử lý nhiều dạng câu hỏi. Tuy nhiên, nếu chỉ dựa
vào tri thức tham số của mô hình, hệ thống có thể gặp các vấn đề:

\begin{itemize}
    \item thông tin học vụ có thể thay đổi theo văn bản hoặc theo học kỳ;
    \item mô hình không biết trạng thái hiện tại của kho dữ liệu nội bộ;
    \item câu trả lời khó kiểm chứng nếu không đi kèm nguồn;
    \item mô hình có thể tạo nội dung nghe hợp lý nhưng không có căn cứ trong tài
    liệu truy hồi.
\end{itemize}

Do đó, LLM trong đồ án đóng vai trò sinh và tổng hợp câu trả lời từ context đã
được truy hồi, thay vì thay thế hoàn toàn kho tri thức.

\subsection{Nguyên lý của RAG}

RAG kết hợp hai quá trình:

\begin{enumerate}
    \item \textbf{Retrieval:} chọn các đoạn dữ liệu có liên quan đến truy vấn từ
    kho đã lập chỉ mục;
    \item \textbf{Augmented Generation:} đưa các đoạn đó vào prompt để LLM sinh
    câu trả lời có căn cứ.
\end{enumerate}

Một vòng đời RAG đầy đủ thường gồm ba pha:

\begin{itemize}
    \item \textbf{Indexing:} chuyển đổi tài liệu, chia chunk, tạo embedding, lưu
    vector và metadata vào các kho truy hồi;
    \item \textbf{Retrieval:} hiểu truy vấn, chọn miền dữ liệu, áp dụng filter,
    tìm candidate và xếp hạng candidate;
    \item \textbf{Generation:} định dạng context, sinh câu trả lời, trả nguồn và
    metadata cần thiết cho giao diện hoặc hệ thống đánh giá.
\end{itemize}

RAG có ưu điểm là cập nhật tri thức bằng cách cập nhật kho dữ liệu thay vì phải
huấn luyện lại LLM sau mỗi lần thay đổi văn bản. Tuy nhiên, RAG không tự bảo đảm
đúng nếu bước chunking, routing hoặc retrieval đưa vào context sai. Vì vậy chất
lượng truy hồi là điều kiện quan trọng của chất lượng câu trả lời.

\section{Xử lý tài liệu và chunking}

\subsection{Chuyển đổi tài liệu sang dạng có thể truy hồi}

Tài liệu đầu vào của hệ thống có thể là PDF, Markdown, HTML đã làm sạch hoặc dữ
liệu dạng JSON từ crawler. Với PDF, bước chuyển đổi sang text hoặc Markdown cần
giữ càng nhiều cấu trúc càng tốt, đặc biệt là heading, bảng, danh sách và thứ bậc
điều khoản.

PDF tiếng Việt có thể phát sinh các vấn đề như font encoding không đồng nhất,
header/footer lặp lại, bảng dài, cột nhiều vùng, scan hoặc ký tự rác do bước
trích xuất. Vì vậy quá trình convert và clean ảnh hưởng trực tiếp đến chất lượng
chunking và retrieval. Trong mã nguồn hiện tại, luồng admin upload dùng các
converter như \texttt{pymupdf4llm} và \texttt{docling}; các pipeline ngoại tuyến
và chunker pháp lý còn hỗ trợ những dữ liệu đầu ra kiểu OCR cho tài liệu quy
định.

\subsection{Khái niệm chunk}

Chunk là đơn vị nội dung được lưu và truy hồi trong hệ thống RAG. Nếu chunk quá
ngắn, nội dung mất ngữ cảnh; nếu chunk quá dài, đoạn văn chứa nhiều thông tin
không liên quan và làm giảm độ chính xác khi đưa vào prompt. Vì vậy một chunk tốt
cần cân bằng:

\begin{itemize}
    \item phạm vi ngữ nghĩa đủ hoàn chỉnh;
    \item kích thước phù hợp với embedding, reranker và ngân sách context;
    \item metadata đủ để xác định nguồn, cấu trúc, collection và phạm vi áp dụng;
    \item khả năng đứng độc lập khi được truy hồi.
\end{itemize}

\subsection{Chiến lược chunking theo loại dữ liệu}

Không có một chiến lược chia đoạn tối ưu cho mọi loại tài liệu. Hệ thống hiện
tại sử dụng các nhóm chiến lược khác nhau:

\begin{table}[h]
\centering
\begin{tabular}{|l|p{4.0cm}|p{6.4cm}|}
\hline
\textbf{Chiến lược} & \textbf{Loại dữ liệu phù hợp} & \textbf{Ý tưởng chính} \\
\hline
Recursive & Markdown hoặc text tổng quát & Tách theo heading, đoạn, dòng và
ranh giới nhỏ hơn khi cần. \\
\hline
Hierarchical/legal & Văn bản quy định có cấu trúc điều khoản & Bảo toàn cấu trúc
chương, điều, khoản, điểm và hierarchy metadata. \\
\hline
\texttt{kehoach} & Kế hoạch và thông báo đã crawl & Tách theo nội dung bài viết,
nguồn và ngày để phục vụ freshness retrieval. \\
\hline
\texttt{stsv} & Sổ tay và dữ liệu hỗ trợ sinh viên & Tách theo mục hướng dẫn,
ngữ cảnh section và cấu trúc dữ liệu nguồn. \\
\hline
\end{tabular}
\caption{Nhóm chiến lược chunking theo dữ liệu}
\end{table}

\subsection{Chunking phân cấp cho văn bản quy định}

Văn bản quy định thường có cấu trúc phân cấp tự nhiên:

\begin{verbatim}
CHUONG I - QUY DINH CHUNG
    Dieu 1. Pham vi dieu chinh
        Khoan 1 ...
        Khoan 2 ...
            Diem a ...
            Diem b ...
\end{verbatim}

Với cấu trúc này, phần cha có thể giữ ngữ cảnh rộng hơn, trong khi phần con phù
hợp cho truy hồi chính xác hơn. Metadata hierarchy giúp hệ thống biết chunk thuộc
điều nào, khoản nào và tài liệu nào. Trong luồng quản trị tài liệu hiện tại,
parent hoặc header chunk vẫn có thể được lưu để review nhưng chính sách indexing
có thể loại chúng khỏi tập truy hồi để tránh chiếm slot ngữ cảnh thay cho child
chunk có nội dung trả lời trực tiếp.

\subsection{Metadata và lineage}

Metadata là cầu nối giữa nội dung văn bản và truy vấn có cấu trúc. Các trường như
\texttt{major\_code}, \texttt{major\_name}, \texttt{applicable\_cohort},
\texttt{applicable\_major}, \texttt{document\_type}, \texttt{date\_str},
\texttt{source} hoặc \texttt{document\_id} giúp hệ thống lọc và giải thích kết
quả truy hồi.

Bên cạnh metadata của từng chunk, lineage của văn bản giúp biểu diễn quan hệ văn
bản mới thay thế văn bản cũ. Với miền quy định, đây là cơ sở để validity filtering
giảm nguy cơ đưa văn bản hết hiệu lực vào context khi đã có thông tin thay thế.

\section{Biểu diễn ngữ nghĩa bằng vector embedding}

\subsection{Khái niệm vector embedding}

Vector embedding ánh xạ văn bản sang một vector số trong không gian nhiều chiều.
Các câu hoặc đoạn có nghĩa gần nhau thường có vector gần nhau theo một độ đo như
cosine similarity. Với hai vector $u$ và $v$, cosine similarity được viết:

\begin{equation}
\operatorname{cos}(u,v) =
\frac{u \cdot v}{\|u\| \|v\|}
\end{equation}

Trong dense retrieval, tài liệu được embed trước ở giai đoạn indexing; truy vấn
được embed tại runtime rồi so khớp với vector tài liệu để tìm các candidate có
ngữ nghĩa gần nhất.

\subsection{BGE-M3 và E5 multilingual}

Hệ thống sử dụng hai mô hình embedding đa ngôn ngữ:

\begin{itemize}
    \item \textbf{BGE-M3:} cung cấp tín hiệu dense retrieval và đồng thời được
    dùng làm đặc trưng embedding cho bộ phân loại miền truy vấn;
    \item \textbf{E5 multilingual:} cung cấp tín hiệu dense retrieval thứ hai,
    phù hợp cho văn bản và truy vấn đa ngôn ngữ.
\end{itemize}

Hai mô hình trong runtime hiện tại tạo hai vector riêng có kích thước 1024 và
được lưu thành named vectors \texttt{bge\_m3} và \texttt{e5} trong Qdrant. Điều
này khác với cách gộp vector bằng trung bình trọng số. Mã nguồn vẫn có
\texttt{EnsembleEmbedder} cho tương thích hoặc thí nghiệm, nhưng đường truy hồi
chính sử dụng hai tín hiệu BGE và E5 riêng rồi hợp nhất ở tầng search.

\subsection{Qdrant và vector store}

Vector store cần hỗ trợ tìm kiếm gần đúng hoặc chính xác theo vector, lưu payload
metadata và lọc theo điều kiện. Trong đồ án, Qdrant được dùng như kho vector theo
collection. Mỗi miền tri thức có collection riêng, vector dùng cosine distance,
và payload chứa text cùng metadata cần thiết cho bước hợp nhất kết quả.

Tách vector store theo miền giúp routing và filtering rõ ràng hơn. Việc dùng named
vectors cho phép một chunk mang đồng thời biểu diễn BGE-M3 và E5 mà không buộc hệ
thống phải mất thông tin của từng tín hiệu.

\section{Tìm kiếm từ khóa và BM25}

\subsection{Vai trò của sparse retrieval}

Dense retrieval mạnh ở ngữ nghĩa, nhưng các câu hỏi học vụ thường chứa token mà
không nên làm mờ: mã môn, mã ngành, số điều khoản, tên biểu mẫu hoặc cụm thuật
ngữ chính xác. Sparse retrieval giữ lợi thế trong những trường hợp đó vì xếp hạng
dựa trên sự xuất hiện của từ khóa trong tài liệu.

Elasticsearch trong hệ thống đảm nhiệm keyword search và metadata search. Ngoài
việc trả kết quả BM25, Elasticsearch còn hỗ trợ bước tìm id theo metadata để áp
điều kiện tương ứng cho vector search.

\subsection{Công thức BM25}

Với truy vấn $Q = \{q_1, q_2, ..., q_n\}$ và tài liệu $D$, điểm BM25 thường được
biểu diễn:

\begin{equation}
\operatorname{BM25}(D,Q) =
\sum_{i=1}^{n}
\operatorname{IDF}(q_i)
\frac{f(q_i,D)(k_1+1)}
{f(q_i,D)+k_1\left(1-b+b\frac{|D|}{avgdl}\right)}
\end{equation}

Trong đó:

\begin{itemize}
    \item $f(q_i,D)$ là tần suất token $q_i$ trong tài liệu $D$;
    \item $|D|$ là độ dài tài liệu;
    \item $avgdl$ là độ dài tài liệu trung bình;
    \item $k_1$ điều khiển độ bão hòa tần suất từ;
    \item $b$ điều khiển mức chuẩn hóa theo độ dài tài liệu.
\end{itemize}

Một dạng IDF thường dùng là:

\begin{equation}
\operatorname{IDF}(q_i) =
\ln\left(
\frac{N-n(q_i)+0.5}{n(q_i)+0.5}+1
\right)
\end{equation}

Trong đó $N$ là số tài liệu trong corpus và $n(q_i)$ là số tài liệu chứa token
$q_i$.

\section{Hybrid retrieval và hợp nhất kết quả}

\subsection{Lý do kết hợp dense và keyword retrieval}

Dense retrieval và BM25 bổ sung cho nhau:

\begin{itemize}
    \item dense retrieval tìm được đoạn diễn đạt khác từ khóa của người dùng;
    \item keyword retrieval giữ tín hiệu chính xác của mã môn, tên riêng và số
    điều khoản;
    \item nếu chỉ dùng một trong hai, hệ thống dễ mất candidate quan trọng trước
    khi reranker được áp dụng.
\end{itemize}

Vì vậy, hybrid retrieval tạo pool candidate từ cả vector search và keyword search
rồi hợp nhất trước khi xếp hạng lại.

\subsection{Score fusion trong tìm kiếm nhiều collection}

Điểm vector và điểm BM25 không cùng thang đo. Một cách hợp nhất là chuẩn hóa từng
pool về đoạn $[0,1]$ rồi cộng theo trọng số:

\begin{equation}
S(d) =
w_v \cdot \widehat{S_v}(d)
+ w_k \cdot \widehat{S_k}(d)
+ B_{\text{recency}}(d)
\end{equation}

Trong đó $\widehat{S_v}$ là điểm vector đã chuẩn hóa, $\widehat{S_k}$ là điểm
keyword đã chuẩn hóa, $w_v$ và $w_k$ là trọng số, còn
$B_{\text{recency}}$ là phần cộng thêm khi collection có nhu cầu ưu tiên độ mới
như \texttt{kehoach}. Theo cấu hình mặc định của tầng
\texttt{MultiCollectionSearch}, vector signal được ưu tiên hơn keyword signal;
với câu hỏi mang tính học phần hoặc mã môn, trọng số keyword có thể được tăng để
giữ exact match.

\subsection{Reciprocal Rank Fusion}

Reciprocal Rank Fusion (RRF) là một kỹ thuật khác để hợp nhất danh sách xếp hạng
khi không muốn phụ thuộc trực tiếp vào thang điểm gốc:

\begin{equation}
\operatorname{RRF}(d) =
\sum_{r \in R}
\frac{1}{k + \operatorname{rank}_r(d)}
\end{equation}

RRF hữu ích khi so sánh nhiều hệ thống retrieval có score khác nhau. Mã nguồn có
hỗ trợ đường fusion theo RRF và một lớp hybrid search theo collection; tuy nhiên
đường tìm kiếm nhiều collection hiện tại mặc định dùng weighted score fusion sau
min--max normalization. Vì vậy trong đồ án, RRF là nền tảng liên quan còn score
fusion là cơ chế cần hiểu khi phân tích runtime chính.

\subsection{Deduplication và candidate pool}

Khi truy hồi từ nhiều collection và nhiều store, cùng một đoạn hoặc các đoạn có
nội dung trùng có thể xuất hiện nhiều lần. Deduplication theo id và theo text giúp
không lãng phí ngân sách context. Candidate pool cũng thường lớn hơn số nguồn trả
về cuối cùng để reranker có đủ lựa chọn trước khi cắt top-$k$.

\section{Metadata-aware retrieval}

\subsection{Lọc theo collection}

Metadata filtering giúp giảm không gian tìm kiếm và tăng xác suất lấy đúng văn
bản. Cơ chế filter phụ thuộc collection:

\begin{itemize}
    \item với \texttt{ctdt}, filter ưu tiên mã ngành hoặc tên ngành;
    \item với \texttt{quydinh}, filter xem xét khóa và ngành áp dụng, đồng thời có
    fallback sang văn bản dùng chung;
    \item với \texttt{kehoach}, ngày đăng bài và yêu cầu freshness có vai trò
    quan trọng;
    \item với \texttt{stsv}, mặc định thường không cần prefilter theo ngành hoặc
    khóa.
\end{itemize}

Một filter tốt không nên làm hệ thống trả rỗng quá sớm. Vì vậy các fallback chain
từ filter chặt sang filter rộng hơn là cần thiết khi metadata không đầy đủ.

\subsection{Phân biệt học kỳ và ngày đăng bài}

Trong dữ liệu học vụ, token như \texttt{2025.2}, \texttt{20252} hoặc
\texttt{2025-2} thường chỉ học kỳ hoặc scope học vụ, không đồng nghĩa với ngày
đăng bài. Việc phân biệt học kỳ với posting date tránh áp nhầm bộ lọc ngày cho
collection \texttt{kehoach}. Với câu hỏi ``mới nhất'' nhưng không chứa ngày cụ
thể, chiến lược phù hợp là ưu tiên nguồn mới thay vì tự suy diễn một học kỳ cũ từ
lịch sử hội thoại.

\subsection{Validity filtering và reference resolution}

Sau retrieval, hai dạng hậu xử lý có ý nghĩa đặc biệt với văn bản quy định:

\begin{itemize}
    \item \textbf{Validity filtering:} loại hoặc hạ ưu tiên tài liệu đã bị thay
    thế khi registry lineage cung cấp thông tin đáng tin cậy;
    \item \textbf{Reference resolution:} phát hiện dẫn chiếu như ``Điều 5'' hoặc
    ``Khoản 1 Điều 5'' và bổ sung chunk liên quan trong cùng văn bản nếu context
    ban đầu chưa đủ.
\end{itemize}

Validity filtering không nên biến câu trả lời thành rỗng khi registry không đủ
thông tin; reference resolution cũng cần ràng buộc theo cùng tài liệu để tránh
ghép điều khoản của văn bản khác.

\section{Reranking bằng cross-encoder}

\subsection{Bi-encoder và cross-encoder}

Embedding retrieval thường hoạt động theo kiểu bi-encoder: query và document
được biểu diễn riêng, sau đó so độ gần trong không gian vector. Cách này nhanh và
phù hợp để tạo candidate pool lớn. Cross-encoder nhận đồng thời cặp
$(query, document)$, cho phép mô hình đánh giá tương tác giữa hai chuỗi chi tiết
hơn nhưng chi phí cao hơn.

\begin{table}[h]
\centering
\begin{tabular}{|l|p{4.2cm}|p{4.2cm}|}
\hline
\textbf{Tiêu chí} & \textbf{Bi-encoder} & \textbf{Cross-encoder} \\
\hline
Đầu vào & Query và document encode tách rời & Cặp query--document encode cùng lúc \\
\hline
Mục đích & Tạo candidate nhanh & Xếp hạng lại candidate \\
\hline
Khả năng mở rộng & Tốt với vector index & Phù hợp cho tập candidate nhỏ hơn \\
\hline
Độ chi tiết so khớp & Phụ thuộc khoảng cách vector & Đánh giá relevance trực tiếp \\
\hline
\end{tabular}
\caption{So sánh bi-encoder và cross-encoder}
\end{table}

\subsection{Vai trò của reranker trong RAG}

Reranker là tầng relevance thứ hai sau hybrid retrieval. Trong hệ thống, BGE
reranker sắp xếp candidate theo điểm liên quan và có thể áp ngưỡng điểm trước khi
cắt top-$k$. Tài liệu bảng, đặc biệt trong chương trình đào tạo, có thể cần ngưỡng
xử lý riêng vì bảng mang nhiều thông tin cấu trúc nhưng văn bản tuyến tính của nó
không luôn được cross-encoder đánh giá giống đoạn prose thông thường.

\section{Hiểu truy vấn và định tuyến}

\subsection{Complexity routing}

Không phải câu hỏi nào cũng cần cùng một luồng xử lý. Complexity routing phân
truy vấn thành:

\begin{itemize}
    \item \texttt{chitchat}: hội thoại ngắn không cần truy hồi;
    \item \texttt{simple}: câu hỏi đủ rõ để đi qua RAG tuyến tính;
    \item \texttt{complex}: câu hỏi cần so sánh, nhiều nguồn, kiểm tra ngữ cảnh cá
    nhân hoặc suy luận nhiều bước.
\end{itemize}

Với nhóm phức tạp, subtype như \texttt{comparison}, \texttt{multi\_source},
\texttt{personal\_check} hoặc \texttt{general} giúp pipeline chọn nhánh phân rã
truy vấn hoặc agent phù hợp.

\subsection{Domain classification và collection selection}

Sau complexity routing, hệ thống cần xác định ý định và miền dữ liệu. Bộ phân
loại miền hiện tại có hai tầng:

\begin{enumerate}
    \item phân loại intent thành \texttt{chitchat}, \texttt{rag} hoặc
    \texttt{tool\_search};
    \item nếu intent là RAG, dự đoán một hoặc nhiều domain trong
    \texttt{ctdt}, \texttt{quydinh}, \texttt{kehoach}, \texttt{stsv}.
\end{enumerate}

Về mặt lý thuyết, domain classification làm giảm nhiễu retrieval và giúp dành
ngân sách candidate cho đúng collection. Trong code hiện tại, embedding BGE-M3
được dùng làm đặc trưng cho classifier học máy; khi độ tin cậy yếu, pipeline có
thể dùng tầng LLM fallback để hiệu chỉnh định tuyến. Với follow-up ngắn và độ tin
cậy thấp, history có thể hỗ trợ route; với câu hỏi dài và tự chứa, history không
nên làm lệch miền dữ liệu.

\subsection{Query reflection}

Query reflection là quá trình biến câu hỏi runtime thành truy vấn truy hồi rõ
ràng hơn. Một reflector có thể:

\begin{itemize}
    \item loại bớt nhiễu hội thoại và thông tin định danh không cần cho retrieval;
    \item viết lại câu hỏi thành dạng standalone khi có đại từ hoặc tham chiếu;
    \item hợp nhất hồ sơ người dùng chỉ khi câu hỏi thực sự phụ thuộc vào hồ sơ;
    \item trích xuất entity như mã ngành, khóa, mã học phần, học kỳ và năm học;
    \item áp guardrail để tránh sinh thêm phạm vi ngành, khóa hoặc học kỳ không
    được câu hỏi hiện tại yêu cầu.
\end{itemize}

Điểm quan trọng là reflection không đồng nghĩa với thêm càng nhiều context càng
tốt. Với câu hỏi mới nhất hoặc câu hỏi chung, lịch sử hội thoại cũ có thể gây
``context bleed'' và làm filter sai. Vì vậy reflection cần có cả khả năng mở rộng
truy vấn và khả năng từ chối context không phù hợp.

\subsection{Query decomposition}

Một câu hỏi có thể bao gồm nhiều mục tiêu retrieval. Query decomposition tách câu
hỏi như vậy thành một số sub-query gắn với collection thích hợp, sau đó hợp nhất
ngữ cảnh trước khi generation. Cách làm này phù hợp với câu hỏi nhiều nguồn hoặc
so sánh vì mỗi truy vấn con có thể tập trung vào một miền dữ liệu rõ ràng.

\section{Agentic RAG và công cụ ngoại vi}

\subsection{Khi nào RAG tuyến tính chưa đủ}

Luồng RAG tuyến tính phù hợp khi một lượt retrieval và một lượt generation đủ để
trả lời. Với câu hỏi cần so sánh, lập kế hoạch tìm kiếm nhiều bước, hoặc quyết
định có cần hỏi lại người dùng hay không, hệ thống cần khả năng điều phối linh
hoạt hơn.

\subsection{Planner--executor và ReAct}

Agentic RAG kết hợp LLM với trạng thái và công cụ. Trong hệ thống hiện tại có hai
hướng xử lý chính:

\begin{itemize}
    \item \textbf{Planner--executor:} dùng cho comparison hoặc multi-source; agent
    phân rã câu hỏi, tạo kế hoạch retrieval có cấu trúc, thực thi các bước tìm
    kiếm và tổng hợp kết quả;
    \item \textbf{ReAct tool loop:} dùng cho câu hỏi phức tạp tổng quát hoặc khi
    planner fallback; LLM chọn công cụ, nhận kết quả, tiếp tục hoặc kết thúc.
\end{itemize}

Các công cụ LLM-visible quan trọng gồm \texttt{rag\_search},
\texttt{web\_search} và \texttt{clarify\_question}. Việc giới hạn tool schema
giúp hành vi agent dễ kiểm soát hơn so với việc cho mô hình gọi tùy ý mọi hàm
trong hệ thống.

\subsection{Web fallback có kiểm soát}

Kho dữ liệu nội bộ phù hợp với quy định và tài liệu đã index, nhưng thông báo có
tính thời điểm có thể cần thông tin mới. Web fallback là cơ chế mở rộng retrieval
ra nguồn ngoài khi:

\begin{itemize}
    \item local retrieval không có nguồn hoặc có confidence thấp;
    \item câu hỏi mang tín hiệu dynamic hoặc freshness;
    \item self-evaluation hoặc agent xác định cần thêm nguồn web.
\end{itemize}

Trong hệ thống, web search là thành phần tùy chọn, dùng danh sách domain cho phép
và ưu tiên nguồn HUST hoặc nguồn giáo dục có thẩm quyền. Đường streaming không
thực hiện hậu kiểm web fallback giống non-streaming path để giữ semantics phát
token theo thời gian thực.

\section{Sinh câu trả lời và tự đánh giá}

\subsection{Prompt grounding}

Sau retrieval, context cần được định dạng cho LLM theo cách giữ được nội dung,
nguồn và các metadata hữu ích. Prompt grounding yêu cầu mô hình trả lời dựa trên
context được cung cấp và thừa nhận thiếu thông tin khi context không đủ. Đây là
cơ chế quan trọng để giảm câu trả lời suy đoán.

Lớp \texttt{llm} trong hệ thống trừu tượng hóa provider qua một interface chung
cho generation thường và streaming. Nhờ đó pipeline không cần gắn chặt với một
SDK cụ thể, trong khi cấu hình runtime vẫn có thể chọn provider phù hợp.

\subsection{Self-evaluation}

Self-evaluation là bước đánh giá phản hồi đã sinh bằng các chiều như:

\begin{itemize}
    \item \textbf{Relevance:} câu trả lời có giải quyết câu hỏi không;
    \item \textbf{Faithfulness:} nội dung có bám vào context không;
    \item \textbf{Completeness:} câu trả lời đã bao quát đủ ý cần thiết chưa.
\end{itemize}

Kết quả self-evaluation có thể hỗ trợ quyết định fallback hoặc đánh dấu answer
status. Tuy nhiên, tự đánh giá bằng LLM không thay thế benchmark retrieval, test
contract và kiểm tra hồi quy trên dữ liệu đã gắn nhãn.

\section{Truy vết, lưu trữ và đánh giá}

\subsection{Truy vết và khả năng quan sát}

Một pipeline RAG có nhiều bước nên cần quan sát được thời gian và quyết định của
từng bước. Trace metadata có thể ghi routing result, reflected question, filter
đã áp dụng, số candidate theo collection, timings, tool calls và agent trace.
Nhờ đó người phát triển phân biệt được lỗi do route sai, filter quá chặt, retrieval
thiếu nguồn hay generation xử lý context chưa tốt.

\subsection{Đánh giá retrieval}

Đánh giá retrieval đo khả năng tìm đúng nguồn trước khi đánh giá câu trả lời cuối
cùng. Các chỉ số thường gặp gồm:

\begin{itemize}
    \item \textbf{Recall@K:} trong top-$K$ có bao nhiêu nguồn liên quan được tìm
    thấy;
    \item \textbf{MRR@K:} vị trí nguồn đúng đầu tiên có cao không;
    \item \textbf{NDCG@K:} xếp hạng có đặt nguồn quan trọng lên gần đầu không;
    \item \textbf{Context precision/recall:} context đưa vào generation vừa đúng
    vừa đủ đến mức nào.
\end{itemize}

Với dữ liệu có tính hiệu lực, cần bổ sung đánh giá freshness và citation/source
accuracy để tránh chỉ tối ưu độ giống từ khóa.

\subsection{Đánh giá hội thoại và đánh giá chính sách hiện hành}

Không nên dùng cùng một bộ kiểm tra cho mọi mục tiêu. Hệ thống hiện tại tách hai
lớp đánh giá:

\begin{itemize}
    \item \textbf{Current policy eval:} kiểm tra retrieval và nguồn factual theo
    dữ liệu hiện đang index;
    \item \textbf{Historical conversation eval:} kiểm tra hiểu follow-up, ngữ cảnh
    cá nhân, chất lượng tư vấn và khả năng làm rõ trong hội thoại lịch sử.
\end{itemize}

Sự tách biệt này quan trọng vì câu trả lời theo email hoặc dữ liệu lịch sử có thể
khác văn bản hiện hành. Một mismatch lịch sử không đồng nghĩa retrieval hiện tại
sai nếu chính sách đã thay đổi.

\section{Kết luận chương}

Chương này đã hệ thống hóa các nền tảng lý thuyết cần thiết cho trợ lý học vụ
RAG trong đồ án:

\begin{itemize}
    \item đặc thù dữ liệu đa miền và yêu cầu trả lời có căn cứ;
    \item xử lý tài liệu, chunking, metadata và lineage;
    \item dense retrieval với BGE-M3/E5, vector store Qdrant và sparse retrieval
    với BM25 trên Elasticsearch;
    \item hybrid retrieval, score fusion, metadata filtering, reranking, validity
    filtering và reference resolution;
    \item query routing, domain classification, reflection, entity extraction và
    decomposition;
    \item Agentic RAG, web fallback có kiểm soát, grounded generation,
    self-evaluation, trace và đánh giá hồi quy.
\end{itemize}

Các cơ sở này là tiền đề để chương tiếp theo trình bày kiến trúc và phương pháp
đề xuất của hệ thống \texttt{RAG\_v2} một cách cụ thể hơn.

\end{document}
